\begin{abstract}
A recent study finds that machine learning can predict the direction of 71\%
of the trades mutual fund managers make, and the trades it cannot predict perform better. What about
hedge funds? We use free SEC filings from 2013 to 2026 and a single modeling
pipeline to study two samples of hedge funds (one hand-picked and one selected
by rule from Form ADV filings made before the sample begins), active long-only
managers, and N-PORT mutual funds. For each existing position in the 500 most
widely held U.S.\ stocks, we predict the quarter-end buy, hold, or sell action
before the filing becomes public. For hedge funds, a gradient-boosted model
corrects only \errredM\% of the mistakes of each fund's majority action
(\gainM{} points of accuracy, in all \dmposM{} test quarters), and none for the median fund. The pipeline finds
planted signals as small as the real effect, and adding more training history provides
little benefit. Hedge funds do not appear to be special: long-only managers and
individual mutual funds are not significantly more predictable, aggregating
mutual funds has no significant effect, and 2017--2019 is no more predictable
than 2022--2026. Three zero-shot time-series foundation models match or fall below the
majority-action baseline in per-fund accuracy. Stocks that the
hand-picked hedge funds sold unexpectedly (mostly against their usual direction) fell more before disclosure than
stocks they sold as predicted, but this effect weakens with size and momentum
matching and does not replicate in the other samples. Apart from this,
unpredicted trades earn no extra return, although with \rqthreequartersM{}
quarters we cannot rule out effects as large as those reported for mutual
funds. Adding \firmn{} firm characteristics barely helps: with public data,
disclosed holdings reveal little about managers' future trades beyond their
habits.
\end{abstract}
